\documentclass[12pt]{article}
\usepackage[utf8]{inputenc}
\usepackage[T1]{fontenc}
\usepackage{hyperref}
\title{Dhuwaya Language Resources for K-12 Curriculum Development}
\author{Compiled by Claude Sonnet 5.5 (claude-sonnet-5-5) for Andr\'as Kornai}
\date{}
\begin{document}
\maketitle

\section{Introduction}
Dhuwaya (ISO 639-3: dwy) is a Yolngu Matha variety of north-east Arnhem Land, Australia. Search results associate it with the long-running bilingual programme at Yirrkala School. Details below come from news and institutional pages, not from the materials themselves.

\section{Primary Educational Resources}
\begin{itemize}
\item Report on Dhuwaya readers at Yirrkala (Indigenous Literacy Foundation): \url{https://www.ilf.org.au/news/dhuwaya-readers-in-yirrkala}. The search summary says 106 readers translated from Djambarrpuyngu into Dhuwaya were produced in 2022, replacing 1980s readers.
\item University of Melbourne news on an early-literacy phonics app teaching in a Yolngu Matha language (reported as Dhuwaya): \url{https://www.unimelb.edu.au/newsroom/news/2022/september/app-teaches-indigenous-early-literacy-for-east-arnhem-land-children-in-yolu-matha-language}
\item Indigenous.gov.au story on teaching mathematics through Yolngu language: \url{https://www.indigenous.gov.au/stories/revolutionising-maths-education-through-yolngu-language}
\end{itemize}

\section{Linguistic and Reference Materials}
None specific to Dhuwaya were found. A scholarly chapter on Yirrkala School's bilingual programme: \url{https://www.doi.org/10.1057/978-1-137-60120-9_4}

\section{Religious and Community Resources}
None identified in the searches.

\section{Digital and Audio Resources}
See the phonics app above.

\section{Curriculum Development Context}
The reference description says Dhuwaya is a language of instruction. Yirrkala School reportedly has a Literature Production Centre and a long-standing bilingual programme; see the NT Department of Education item on the Yirrkala language centre: \url{https://education.nt.gov.au/statistics-research-and-strategies/indigenous-education-strategy/indigenous-education-strategy-issue-17/yirrkala-language-centre-recognised}. Access to the actual materials would need to go through the school and community.
\end{document}